\documentclass[11pt,a4paper]{article}
\usepackage{fontspec,amsmath,amssymb,amsthm,microtype}
\setmainfont{lmroman10-regular.otf}[BoldFont=lmroman10-bold.otf,ItalicFont=lmroman10-italic.otf,BoldItalicFont=lmroman10-bolditalic.otf]
\usepackage{xcolor}
\usepackage[most]{tcolorbox}
\usepackage{mathrsfs}
\providecommand{\lt}{<}
\providecommand{\gt}{>}
\setmonofont{DejaVuSansMono.ttf}[BoldFont=DejaVuSansMono-Bold.ttf,ItalicFont=DejaVuSansMono-Oblique.ttf,BoldItalicFont=DejaVuSansMono-BoldOblique.ttf,Scale=0.86]
\newfontfamily\leanfallback{STIXTwoMath-Regular.otf}[Scale=0.86]
\definecolor{leanrule}{gray}{0.62}
\definecolor{leanback}{gray}{0.962}
\definecolor{leantitle}{gray}{0.22}
\definecolor{leanmuted}{gray}{0.38}
\newtcolorbox{leanbox}[1]{enhanced,breakable,colback=leanback,colframe=leanrule,
  boxrule=0.4pt,arc=0.6mm,left=2.4mm,right=2.4mm,top=1.8mm,bottom=1.8mm,
  before skip=0.8em,after skip=1em,fonttitle=\small\bfseries,coltitle=leantitle,
  colbacktitle=leanback,titlerule=0.3pt,title={#1}}
\newlength{\leancharwidth}
\newenvironment{leancode}{\par\smallskip\begingroup\ttfamily\footnotesize\raggedright
  \settowidth{\leancharwidth}{M}\setlength{\parskip}{0pt}\setlength{\parindent}{0pt}
  \frenchspacing}{\par\endgroup\smallskip}
\newcommand{\leanline}[2]{\hangindent=\dimexpr#1\leancharwidth+4\leancharwidth\relax
  \hangafter=1\noindent\hspace*{#1\leancharwidth}#2\par}
\newcommand{\leansource}[1]{\par{\footnotesize\color{leanmuted}#1}\par}
\newcommand{\leanlabel}[1]{\par\smallskip\noindent{\small\bfseries #1}\par}
\usepackage[margin=25mm]{geometry}
\usepackage[unicode,hidelinks]{hyperref}
\hypersetup{pdfauthor={LangevinLean Project}}
\setlength{\emergencystretch}{3em}
\setlength{\parskip}{0.4em}
\setlength{\parindent}{0pt}
\allowdisplaybreaks[2]
\raggedbottom
\title{\texorpdfstring{Strong convexity of the Moreau envelope of a strongly convex function}{Strong convexity of the Moreau envelope of a strongly convex function}}
\author{LangevinLean Project}
\date{}
\begin{document}\maketitle
Draft mathematical exposition; these notes may contain mistakes.


Let \(E\) be a finite-dimensional real inner-product space, \(m\ge0\), \(\lambda>0\), and let \(f:E\to\mathbb R\) be \(m\)-strongly convex and bounded below. This section proves that the Moreau envelope \(f^\lambda(x)=\min_p\bigl(f(p)+\|p-x\|^2/(2\lambda)\bigr)\) is \(m/(1+m\lambda)\)-strongly convex. The proof tests the minimization that defines \(f^\lambda(ax+by)\) with the convex combination \(aP_\lambda(x)+bP_\lambda(y)\) of the two proximal points, and combines the strong-convexity deficit of \(f\) with the deficit of the quadratic penalty by completing a square. Differentiability of \(f\) is not used.




The result is used when the whole strongly convex potential is smoothed, as in PIPULA. When only a convex summand \(g_2\) is smoothed and an \(m\)-strongly convex \(g_1\) is kept, the sum \(g_1+g_2^\lambda\) keeps the modulus \(m\), because \(g_2^\lambda\) is convex.


\section{\texorpdfstring{Two quadratic identities}{Two quadratic identities}}


\textbf{Claim.} Let \(E\) be a real inner-product space and \(u,v\in E\).



\begin{enumerate}


\item If \(a,b\ge0\) and \(a+b=1\), then
\[\|au+bv\|^2=a\|u\|^2+b\|v\|^2-ab\|u-v\|^2.\qquad(1)\]

\item If \(m\ge0\) and \(\lambda>0\), then
\[\frac{m}{1+m\lambda}\|v\|^2\le m\|u\|^2+\frac{\|u-v\|^2}{\lambda}.\qquad(2)\]

\end{enumerate}



\textbf{Proof.} (a) Expanding the squared norms by bilinearity of the inner product,


\[\|au+bv\|^2=a^2\|u\|^2+2ab\langle u,v\rangle+b^2\|v\|^2,\qquad ab\|u-v\|^2=ab\|u\|^2-2ab\langle u,v\rangle+ab\|v\|^2.\]


Hence \(a\|u\|^2+b\|v\|^2-ab\|u-v\|^2=a(1-b)\|u\|^2+2ab\langle u,v\rangle+b(1-a)\|v\|^2\), and since \(1-b=a\) and \(1-a=b\), this equals \(\|au+bv\|^2\).




(b) Put \(c=1+m\lambda\). Since \(m\ge0\) and \(\lambda>0\), \(c\ge1>0\). We show


\[m\|u\|^2+\frac{\|u-v\|^2}{\lambda}-\frac{m}{c}\|v\|^2=\frac{\|cu-v\|^2}{\lambda c}.\qquad(3)\]


Multiply the left side by \(\lambda c>0\) and expand \(\|u-v\|^2=\|u\|^2-2\langle u,v\rangle+\|v\|^2\):


\[\begin{aligned}
m\lambda c\|u\|^2+c\bigl(\|u\|^2-2\langle u,v\rangle+\|v\|^2\bigr)-m\lambda\|v\|^2
&=c(1+m\lambda)\|u\|^2-2c\langle u,v\rangle+(c-m\lambda)\|v\|^2\\
&=c^2\|u\|^2-2c\langle u,v\rangle+\|v\|^2=\|cu-v\|^2 .
\end{aligned}\]


Dividing by \(\lambda c\) gives (3). The right-hand side of (3) is nonnegative, because it is a square divided by the positive number \(\lambda c\); this is (2). \(\square\)




For \(v=cu\) the right-hand side of (3) vanishes, so equality holds in (2) and the constant \(m/(1+m\lambda)\) cannot be replaced by a larger one.

\begin{leanbox}{Lean: \texttt{norm\_\allowbreak{}convexCombo\_\allowbreak{}sq}, \texttt{weighted\_\allowbreak{}quadratic\_\allowbreak{}lower}}
\leansource{\texttt{norm\_\allowbreak{}convexCombo\_\allowbreak{}sq} \quad ProximalLangevin/\allowbreak{}StrongMoreau.lean, lines 21–29}
\begin{leancode}
\leanline{0}{theorem\ norm\_convexCombo\_sq\ (u\ v\ :\ E)\ \{a\ b\ :\ ℝ\}\ (ha\ :\ 0\ ≤\ a)\ (hb\ :\ 0\ ≤\ b)}
\leanline{4}{(hab\ :\ a\ +\ b\ =\ 1)\ :}
\leanline{4}{‖a\ •\ u\ +\ b\ •\ v‖\ \textasciicircum{}\ 2\ =\ a\ *\ ‖u‖\ \textasciicircum{}\ 2\ +\ b\ *\ ‖v‖\ \textasciicircum{}\ 2\ -}
\leanline{6}{a\ *\ b\ *\ ‖u\ -\ v‖\ \textasciicircum{}\ 2}
\end{leancode}
\leansource{\texttt{weighted\_\allowbreak{}quadratic\_\allowbreak{}lower} \quad ProximalLangevin/\allowbreak{}StrongMoreau.lean, lines 32–46}
\begin{leancode}
\leanline{0}{theorem\ weighted\_quadratic\_lower\ \{m\ lam\ :\ ℝ\}\ (hm\ :\ 0\ ≤\ m)\ (hlam\ :\ 0\ <\ lam)}
\leanline{4}{(u\ v\ :\ E)\ :}
\leanline{4}{m\ /\ (1\ +\ m\ *\ lam)\ *\ ‖v‖\ \textasciicircum{}\ 2\ ≤\ m\ *\ ‖u‖\ \textasciicircum{}\ 2\ +\ ‖u\ -\ v‖\ \textasciicircum{}\ 2\ /\ lam}
\end{leancode}
\leanlabel{How the Lean proof follows the argument}
\texttt{norm\_\allowbreak{}convex\allowbreak{}Combo\_\allowbreak{}sq} is (1) with the hypotheses \texttt{ha : 0 \ensuremath{\le} a}, \texttt{hb : 0 \ensuremath{\le} b} and \texttt{hab : a + b = 1}. The proof rewrites with \texttt{norm\_\allowbreak{}add\_\allowbreak{}sq\_\allowbreak{}real}, \texttt{norm\_\allowbreak{}smul}, \texttt{Real.\allowbreak{}norm\_\allowbreak{}eq\_\allowbreak{}abs} and \texttt{abs\_\allowbreak{}of\_\allowbreak{}nonneg} (here the signs of \(a\) and \(b\) are used to remove \(|a|\) and \(|b|\)), \texttt{mul\_\allowbreak{}pow}, \texttt{real\_\allowbreak{}inner\_\allowbreak{}smul\_\allowbreak{}left}, \texttt{real\_\allowbreak{}inner\_\allowbreak{}smul\_\allowbreak{}right} and \texttt{norm\_\allowbreak{}sub\_\allowbreak{}sq\_\allowbreak{}real}; then \texttt{obtain rfl := eq\_\allowbreak{}sub\_\allowbreak{}of\_\allowbreak{}add\_\allowbreak{}eq hab} substitutes \(a=1-b\) and \texttt{ring} verifies the resulting polynomial identity.




\texttt{weighted\_\allowbreak{}quadratic\_\allowbreak{}lower} is (2) under \texttt{hm : 0 \ensuremath{\le} m} and \texttt{hlam : 0 < lam}. \texttt{hden} is \(c>0\); \texttt{heq} is identity (3), proved by expanding both squared norms, extracting the factor \(c\) with \texttt{norm\_\allowbreak{}smul} and \texttt{abs\_\allowbreak{}of\_\allowbreak{}pos hden}, and finishing with \texttt{field\_\allowbreak{}simp} and \texttt{ring}; \texttt{hpos} is the nonnegativity of the right side of (3) (\texttt{div\_\allowbreak{}nonneg}); rewriting \texttt{hpos} with \texttt{← heq} and \texttt{linarith} give (2).
\end{leanbox}

\section{\texorpdfstring{The envelope of an \(m\)-strongly convex function is \(m/(1+m\lambda)\)-strongly convex}{The envelope of an m-strongly convex function is m/(1+mλ)-strongly convex}}


Recall that \(f:E\to\mathbb R\) is \emph{\(m\)-strongly convex} if for all \(x,y\in E\) and all \(a,b\ge0\) with \(a+b=1\),


\[f(ax+by)\le af(x)+bf(y)-ab\,\frac m2\,\|x-y\|^2.\]


For \(m=0\) this is convexity; for \(m\ge0\) an \(m\)-strongly convex function is convex.




\textbf{Claim.} Let \(E\) be finite-dimensional, \(m\ge0\), \(\lambda>0\), and let \(f:E\to\mathbb R\) be \(m\)-strongly convex and bounded below. Then \(f^\lambda\) is \(m_\lambda\)-strongly convex, where


\[m_\lambda=\frac{m}{1+m\lambda}.\]


\textbf{Proof.} Since \(m\ge0\), \(f\) is convex; together with the lower bound and finite dimension this gives the proximal map \(P_\lambda\) and the envelope \(f^\lambda\) of the previous sections. Let \(x,y\in E\), \(a,b\ge0\) with \(a+b=1\), and put \(z=ax+by\), \(p=P_\lambda(x)\), \(q=P_\lambda(y)\). Since \(P_\lambda(z)\) minimizes \(Q_{\lambda,z}\), the competitor \(ap+bq\) gives


\[f^\lambda(z)\le f(ap+bq)+\frac{\|ap+bq-z\|^2}{2\lambda}.\qquad(4)\]


Since \(ap+bq-z=a(p-x)+b(q-y)\), identity (1) of the previous step with \(u=p-x\) and \(v=q-y\), for which \(u-v=(p-q)-(x-y)\), gives


\[\|ap+bq-z\|^2=a\|p-x\|^2+b\|q-y\|^2-ab\,\|(p-q)-(x-y)\|^2.\]


Strong convexity of \(f\) at \(p,q\) gives \(f(ap+bq)\le af(p)+bf(q)-ab\frac m2\|p-q\|^2\). Inserting both into (4) and regrouping,


\[f^\lambda(z)\le a\Bigl(f(p)+\frac{\|p-x\|^2}{2\lambda}\Bigr)+b\Bigl(f(q)+\frac{\|q-y\|^2}{2\lambda}\Bigr)-\frac{ab}{2}\Bigl(m\|p-q\|^2+\frac{\|(p-q)-(x-y)\|^2}{\lambda}\Bigr).\]


The first two brackets are \(f^\lambda(x)\) and \(f^\lambda(y)\). By (2) of the previous step with \(u=p-q\) and \(v=x-y\), the last bracket is at least \(m_\lambda\|x-y\|^2\); since \(ab\ge0\), the last term is at most \(-\frac{ab}{2}m_\lambda\|x-y\|^2\). Therefore


\[f^\lambda(ax+by)\le af^\lambda(x)+bf^\lambda(y)-ab\,\frac{m_\lambda}{2}\,\|x-y\|^2 .\]


This is \(m_\lambda\)-strong convexity of \(f^\lambda\). \(\square\)




No derivative of \(f\) appears in the argument. For \(m=0\) the claim reduces to convexity of the envelope.

\begin{leanbox}{Lean: \texttt{moreau\_\allowbreak{}strongConvex}}
\leansource{\texttt{moreau\_\allowbreak{}strongConvex} \quad ProximalLangevin/\allowbreak{}StrongMoreau.lean, lines 49–84}
\begin{leancode}
\leanline{0}{theorem\ moreau\_strongConvex\ [FiniteDimensional\ ℝ\ E]\ \{f\ :\ E\ →\ ℝ\}}
\leanline{4}{\{m\ lam\ :\ ℝ\}\ (hm\ :\ 0\ ≤\ m)\ (hsc\ :\ StrongConvexOn\ univ\ m\ f)}
\leanline{4}{(hb\ :\ ∃\ b,\ ∀\ z,\ b\ ≤\ f\ z)\ (hlam\ :\ 0\ <\ lam)\ :}
\leanline{4}{StrongConvexOn\ univ\ (m\ /\ (1\ +\ m\ *\ lam))\ (moreau\ f\ lam)}
\end{leancode}
\leanlabel{How the Lean proof follows the argument}
\texttt{Strong\allowbreak{}Convex\allowbreak{}On univ m f} is Mathlib's strong convexity on the whole space: it unfolds to \texttt{Uniform\allowbreak{}Convex\allowbreak{}On} with modulus \(r\mapsto\frac m2r^2\), a structure whose first component is convexity of the domain and whose second component is the displayed inequality for weights \(a,b\ge0\), \(a+b=1\).




\texttt{moreau\_\allowbreak{}strong\allowbreak{}Convex} assumes \texttt{[Finite\allowbreak{}Dimensional \ensuremath{\mathbb R} E]}, \texttt{hm : 0 \ensuremath{\le} m}, \texttt{hsc : Strong\allowbreak{}Convex\allowbreak{}On univ m f}, the lower bound \texttt{hb} and \texttt{hlam : 0 < lam}, and concludes \texttt{Strong\allowbreak{}Convex\allowbreak{}On univ (m /\allowbreak{} (1 + m * lam)) (moreau f lam)}.



\begin{enumerate}


\item \texttt{hf} is convexity of \(f\): \texttt{hsc.\allowbreak{}mono hm} weakens the modulus from \(m\) to \(0\), and \texttt{strong\allowbreak{}Convex\allowbreak{}On\_\allowbreak{}zero} identifies \(0\)-strong convexity with convexity.

\item \texttt{refine \ensuremath{\langle}convex\_\allowbreak{}univ, fun x \_\allowbreak{} y \_\allowbreak{} a b ha hb' hab => ?\_\allowbreak{}\ensuremath{\rangle}} introduces the two points and the weights. \texttt{hmin} is (4): \texttt{prox\_\allowbreak{}is\allowbreak{}Prox\allowbreak{}Point} at \texttt{a \ensuremath{\bullet} x + b \ensuremath{\bullet} y} tested at \texttt{a \ensuremath{\bullet} p + b \ensuremath{\bullet} q}. \texttt{hs} is the strong convexity of \(f\) at \(p,q\).

\item \texttt{heq} is \(ap+bq-z=a(p-x)+b(q-y)\); the rewrite with \texttt{norm\_\allowbreak{}convex\allowbreak{}Combo\_\allowbreak{}sq (p -\allowbreak{} x) (q -\allowbreak{} y) ha hb' hab} expands \(\|ap+bq-z\|^2\), and \texttt{hd} rewrites \((p-x)-(q-y)\) as \((p-q)-(x-y)\).

\item \texttt{hw} is \texttt{weighted\_\allowbreak{}quadratic\_\allowbreak{}lower} with \(u=p-q\), \(v=x-y\), and \texttt{hww} multiplies it by \(ab\ge0\). \texttt{hscale}, \texttt{hmin'} and \texttt{hweight} multiply the relevant inequalities by \(2\lambda\) or \(\lambda\); \texttt{field\_\allowbreak{}simp} clears the denominators, and \texttt{nlinarith} combines everything as in the final display.

\end{enumerate}
\end{leanbox}

\end{document}
